%!TEX root = ../thesis.tex

\section{軟弱地面における床沈み込み現象}

軟弱地面とは，砂地や土壌，柔軟材料など，外力に応じて変形を生じる床環境を指す．このような環境では，ロボットの足が接地した際に床が鉛直方向に変形し，足先が床内部へ沈み込む現象が生じる．

剛体床においては，足先の接地位置および床反力の作用点は接地後にほぼ不変であるのに対し，軟弱地面では接地後も床変形が進行し，足先位置および床反力の作用条件が時間的に変化する．この床沈み込み量は，床の剛性や減衰特性，ロボットの質量や接地速度などに依存し，歩行中に逐次的に変動する外乱としてロボットの運動に影響を与える．